\documentclass[a4paper,11pt]{article}

\usepackage{revy}
\usepackage[utf8]{inputenc}
\usepackage[T1]{fontenc}
\usepackage[danish]{babel}
\usepackage{amsmath,amssymb,amsfonts,amsthm,textcomp}


\revyname{MatematikRevyen}
\revyyear{2026}
% HUSK AT OPDATERE VERSIONSNUMMER
\version{0.1}
\eta{$3-4$ minutter}
\status{Færdig}

\title{Løvens A-lokale 1}
\author{Carl '21 \& Thais '22}

\begin{document}
\maketitle

\begin{roles}
\role{E}[Skuespiller] Studerende-entrepeneur (selvglad type)
\role{S}[Skuespiller] Speaker (meget få linjer, ku være i bandet?)
\role{L1}[Skuespiller] Forelæser-løve (Fabien)
\role{L2}[Skuespiller] Forelæser-løve (Ernst)
\role{L3}[Skuespiller] Forelæser-løve (Morten)
\end{roles}

\begin{props}
\item Bord med tre stole til løverne
\end{props}


\begin{sketch}
\says{S} Den næste iværksætter er S, der gerne vil hjælpe forelæsere med at tjene lidt ekstre lommepenge

\says{E} Hey løver, prøv lige at dyrk den her. (viser en bog). Woah, det godt nok en fed bog hva?

\says{L1} Jo, den ser da ret sej ud
\says{L2} Den ku jeg godt se mig selv læse i

\says{E} Men løver, så prøv lige at se her. Årgh hvad, det godt nok en rimeligt tjekket. (peger på den første bog). Prøv lige se den der bog der, fuck den er fyldt med fejl mand, helt ubrugelig. Prøv at se en grim bog. Det her det er fandme en god bog

\scene{Løverne er forvirrede}

\says{L3}[peger på første bog] Men sagde du ikke lige før, at den der bog var god?

\says{E} Jo, SIDSTE år. Og det er kernen i mit koncept: Udgaver (*håndtegn*)

\scene{Intro-jingle}

\says{E} Hvis jeg nu sælger dig en dyr bog. Hvor længe kan du så bruge den?

\says{L2} En dyr bog, kan man vel bruge hele sin uddannelse og sælge videre til nye studerende

\says{E} Men hvad nu hvis jeg sælger dig en dyr bog og så året efter, siger at den er fuld af fejl og den ikke længere kan bruges til noget, hvad gør du så?

\says{L2} Så har jeg nok ikke lyst til at bruge den bog

\says{E} Nemlig! Vi skal have fremelsket en kultur hvor ingen kan bestå, hvis ikke de har den udgave af kursusbogen, som vi siger kan bruges. Året efter, så siger vi så, at det er en ny udgave af bogen, der er den rigtige. Og den gamle bog overhovedet ikke kan bruges mere og så skal de købe en ny bog af os. 

\says{L3} Men hvorfor skulle de studerende lytte til hvad du siger?

\says{E} Øhhh, fordi, at det er mit kursus og jeg laver eksamen, ik? 
\says{E} Og så fordi, at der står mit navn på bogen, så kan vi tage 3-400kr for en bog, der har kostet 25kr i print oppe på 4. 

\says{L1} ror du virkelig at folk er dumme nok, til at betale massere penge for den samme bog, bare fordi det er en ny udgave?

\says{E}[peger på L2] Altså, du kunne da har brug for det, ik? Har du set din bog mand, side 80, mangler et komma i en definition og her har du brugt epsilon i stedet for varepsilon. 

\says{L1}[drillende] Altså, det er lidt pinligt, ik? 
\says{L3}[drillende] Velkommen til dansk.
\says{L2} I sagde da godt I kunne lide min bog:(

\says{E} Vi skal bilde de studerende ind, at de her nye udgave alle sammen er forskellige. Og at de er helt nødvendige for at bestå. De studerende kommer til at æde det råt.

\says{L3} Ej, det er simpelthen for dumt. Hvad så med de overskydende eksemplarer af gamle udgaver, vi har liggende?

\sasy{E} Dem kan vi da også bare sælge! Vi bestemmer. Et år er det udgave 5, der er den rigtige, næste år kan man kun bruge udgave 1-3. 

\says{L1} Altså, det der kommer alle til at gennemskue. Jeg er ude.

\says{E} Jeg siger bare, folk de kommer til at ligger i telt foran bogladen for at få fat i de nye udgaver. 

\says{L3} I telt? Jeg er også ude.

\sasy{L2} Hvis jeg nu går med, kan du så hjælp mig med at rette mine kommafejl?

\sasy{E} Er du gal! Hvis du går med, så retter vi bare de fejl og smider et nyt tal på forsiden! Bum! Og så gør vi det igen næste år!

\says{L2} Er min bog så sej?

\sasy{E} Er du med?! 

\says{L2} Selvfølgelig er jeg med.

\says{E} Fedt! 
\scene{E slapper et nyt tal på deres bog}
\says{E} Holdt kæft en fed udgave er den bog der.

\scene{Lys ned}


\end{sketch}
\end{document}
